% Section 29: Common mistakes to avoid} | Part: 7 | Depends on: see docs/REWRITE_PLAN.md | Status: todo | Original lines: 4084-4097
%!TEX root = ../../../main.tex
\subsection{Common mistakes to avoid}
\begin{enumerate}
\item \textbf{Calling the array $\Gamma^k{}_{ij}$ a tensor.} It is not; $\nabla\tau$ and $R$ are tensors.
\item \textbf{Mixing up covariant and contravariant.} Test the type of an object by its transformation law: if it uses $\partial\tilde x/\partial x$ it is upper-index, if $\partial x/\partial\tilde x$ it is lower-index.
\item \textbf{Treating the gradient as a vector} without a metric ($\dd f$ is a covector; $\nabla f=(\dd f)^\sharp$ needs $g$).
\item \textbf{Assuming $\Gamma=0$ means flat.} Flat means $R=0$.
\item \textbf{Forgetting the Lorentzian signs}: $g(v,v)$ can be negative or zero for $v\neq0$.
\item \textbf{Mixing curvature conventions} between books, \textbf{or signature conventions}: mostly plus versus mostly minus (Section~\ref{sec:conventions}).
\item \textbf{Confusing $\dd x$ with ``a small increment'' and $\partial/\partial x$ with a ratio}: the first is a covector, the second depends on the whole chart (Section~\ref{sec:trio}).
\item \textbf{Thinking Stokes needs a metric.} It needs only a smooth structure and an orientation; the metric enters when forms are turned into vector fields (Section~\ref{sec:classical}).
\item \textbf{Forgetting a factor $k!$} when moving between the two realisations of the wedge (Remark~\ref{rem:kfactorial}).
\item \textbf{Conflating $\nabla$ and $\vd$}, or $R$ and the rotor $\rotor$.
\end{enumerate}
